% ===== PRÉAMBULE CANONIQUE — à coller VERBATIM en tête de CHAQUE .tex =====
\documentclass[11pt,a4paper]{article}
\usepackage[utf8]{inputenc}\usepackage[T1]{fontenc}\usepackage{lmodern}
\usepackage[french]{babel}
\usepackage{amsmath,amssymb,mathtools}\usepackage{icomma}
\usepackage{siunitx}\sisetup{locale=FR}  % \qty{}{}, \si{}, \ang{}
\usepackage{xcolor}\usepackage{colortbl}
\usepackage{tikz}\usepackage{pgfplots}\pgfplotsset{compat=1.18}
\usepackage[siunitx,europeanresistors]{circuitikz} % circuits (élec.)
\usepackage[most]{tcolorbox}\usepackage{enumitem}\usepackage{array,booktabs}
\usepackage[a4paper,margin=2.2cm]{geometry}
\usetikzlibrary{arrows.meta,calc,angles,quotes,positioning,patterns,%
  backgrounds,shapes.geometric,decorations.pathmorphing,decorations.markings,intersections}
\definecolor{primary}{RGB}{222,49,99}\definecolor{primarydark}{RGB}{150,22,55}
\definecolor{defcol}{RGB}{0,114,178}\definecolor{methcol}{RGB}{52,92,148}
\definecolor{excol}{RGB}{0,135,90}\definecolor{attncol}{RGB}{200,40,55}
\tcbset{boxbase/.style={enhanced,breakable,boxrule=0.9pt,arc=2.5pt,
  left=3mm,right=3mm,top=2.3mm,bottom=2.3mm,fonttitle=\bfseries,coltitle=white}}
% --- boîtes TUTORIEL (noms EXACTS, ne pas renommer) ---
\newtcolorbox{cle}[1][]{boxbase,colback=defcol!6,colframe=defcol,
  title={\textsc{À retenir}\ifx\relax#1\relax\relax\else\textmd{ : \itshape #1}\fi}}
\newtcolorbox{methode}[1][]{boxbase,colback=methcol!7,colframe=methcol,sharp corners,
  title={\textsc{Méthode}\ifx\relax#1\relax\relax\else\textmd{ --- \itshape #1}\fi}}
\newtcolorbox{exemplebox}[1][]{boxbase,colback=excol!5,colframe=excol,
  title={\textsc{Exemple}\ifx\relax#1\relax\relax\else\textmd{ : \itshape #1}\fi}}
\newtcolorbox{piege}[1][]{boxbase,colback=attncol!6,colframe=attncol,
  title={\textsc{Piège}\ifx\relax#1\relax\relax\else\textmd{ : \itshape #1}\fi}}
\newtcolorbox{remarque}[1][]{boxbase,colback=black!4,colframe=black!45,
  title={\textsc{Remarque}\ifx\relax#1\relax\relax\else\textmd{ : \itshape #1}\fi}}
% --- boîtes EXERCICE / CORRIGÉ (noms EXACTS, auto-comptées) ---
\newtcolorbox[auto counter]{exo}[1][]{boxbase,colback=primary!4,colframe=primary,
  title={\textsc{Exercice}~\thetcbcounter\ifx\relax#1\relax\relax\else\textmd{ --- \itshape #1}\fi}}
\newtcolorbox[auto counter]{corrige}[1][]{boxbase,colback=excol!4,colframe=excol,
  title={\textsc{Corrigé}~\thetcbcounter\ifx\relax#1\relax\relax\else\textmd{ --- \itshape #1}\fi}}
\newcommand{\vv}[1]{\vec{#1}}\newcommand{\ds}{\displaystyle}
\newcommand{\R}{\mathbb{R}}\newcommand{\N}{\mathbb{N}}\newcommand{\Z}{\mathbb{Z}}
% (un \usepackage{fancyhdr}+en-tête est possible ; il est ignoré côté web)
% =========================================================================
\begin{document}
\sisetup{per-mode=symbol}

\begin{corrige}[intensité de la force de Coulomb]
\begin{enumerate}[label=\alph*)]
  \item Conversions : $q_1 = \num{3e-6}\ \si{\coulomb}$, $q_2 = \num{2e-6}\ \si{\coulomb}$,
        $d = \SI{0.10}{\metre}$.
        \[ F = K\,\frac{|q_1||q_2|}{d^2}
             = \num{9e9}\times\frac{(\num{3e-6})(\num{2e-6})}{(0{,}10)^2}
             = \num{9e9}\times\frac{\num{6e-12}}{0{,}01} = \SI{5.4}{\newton}. \]
  \item Les deux charges sont positives (mêmes signes) : force \textbf{répulsive}.
\end{enumerate}
\end{corrige}

\begin{corrige}[attraction ou répulsion ?]
\begin{enumerate}[label=\alph*)]
  \item Signes opposés ($+$ et $-$) : \textbf{attraction}. La force sur la charge de gauche pointe
        \emph{vers la droite} (vers $q_2$).
  \item Mêmes signes ($-$ et $-$) : \textbf{répulsion}. La force sur la charge de gauche pointe
        \emph{vers la gauche} (elle est repoussée, dos à $q_2$).
\end{enumerate}
\end{corrige}

\begin{corrige}[jouer sur la charge et la distance]
On raisonne par proportionnalité, sans tout recalculer.
\begin{enumerate}[label=\alph*)]
  \item $F \propto 1/d^2$ : doubler $d$ divise $F$ par $2^2 = 4$.
        \[ F = \frac{F_0}{4} = \frac{8}{4} = \SI{2}{\newton}. \]
  \item $F \propto |q|$ : tripler une charge triple $F$.
        \[ F = 3\,F_0 = 3\times 8 = \SI{24}{\newton}. \]
\end{enumerate}
\end{corrige}

\begin{corrige}[retrouver la distance]
\begin{enumerate}[label=\alph*)]
  \item De $F = K\,|q_1||q_2|/d^2$ on isole $d$ :
        \[ d = \sqrt{\frac{K\,|q_1|\,|q_2|}{F}}. \]
  \item \[ d = \sqrt{\frac{\num{9e9}\times(\num{5e-6})(\num{2e-6})}{9}}
             = \sqrt{\frac{\num{9e9}\times\num{1e-11}}{9}} = \sqrt{0{,}01}
             = \SI{0.10}{\metre} = \SI{10}{\centi\metre}. \]
\end{enumerate}
\end{corrige}

\begin{corrige}[la plus grosse charge subit-elle plus ?]
\begin{enumerate}[label=\alph*)]
  \item La distance et le produit des charges sont les mêmes pour les deux, donc une seule valeur :
        \[ F = K\,\frac{|q_1||q_2|}{d^2}
             = \num{9e9}\times\frac{(\num{2e-6})(\num{8e-6})}{(0{,}20)^2}
             = \num{9e9}\times\frac{\num{1.6e-11}}{0{,}04} = \SI{3.6}{\newton}. \]
        $q_1$ subit $\SI{3.6}{\newton}$, $q_2$ subit aussi $\SI{3.6}{\newton}$.
  \item Non : les deux forces sont \textbf{égales} (action--réaction). La formule ne fait
        intervenir que le \emph{produit} $|q_1||q_2|$, identique pour les deux charges. Une charge
        plus grande ne subit pas une force plus intense.
\end{enumerate}
\end{corrige}

\end{document}
